\textbf{Proof.}
Write the four bounded response coordinates as \((a,b,c,d)\) and \(t=(b-a+d-c)/2\). The mapping tables give observations
\[
\begin{array}{c|cc}
&10&01\\ \hline
e^{(0)}&(a,c)&(a,d)\\
e^{(-)}&(a,c)&(b,c)\\
e^{(+)}&(b,c)&(a,c).
\end{array} \tag{1}
\]
Under the uniform design, use respectively
\[
\begin{array}{c|cc}
&10&01\\ \hline
e^{(0)}&-a/2-2c/3&-a/2+2d/3\\
e^{(-)}&-2a/3-c/2&2b/3-c/2\\
e^{(+)}&2b/3-c/2&-2a/3-c/2.
\end{array} \tag{2}
\]
Substitution of the sixteen sign vectors into the two-cell mean squared error gives maximum \(25/36\) in each row. These affine rules are their own multilinear extensions, so \(\mathrm{lem{:}sign\text{-}reduction}\) gives the same upper bound for every bounded response array.

For matching lower bounds, use the four-atom priors
\[
\begin{array}{c|c|c}
e^{(0)}&7/18&(-1,-1,1,-1)\\
&5/18&(-1,1,-1,-1)\\
&1/18&(-1,1,-1,1)\\
&5/18&(-1,1,1,1)\\ \hline
e^{(-)}&5/18&(-1,-1,-1,1)\\
&1/18&(-1,1,-1,1)\\
&7/18&(1,-1,-1,-1)\\
&5/18&(1,1,-1,1)\\ \hline
e^{(+)}&5/18&(-1,-1,-1,-1)\\
&7/18&(-1,1,-1,1)\\
&1/18&(1,-1,-1,-1)\\
&5/18&(1,1,-1,-1).
\end{array} \tag{3}
\]
Direct grouping using (1) gives, at each selector cell and for each mapping, two nonempty posterior cells whose variance contributions are \(35/54\) and \(5/108\). Their sum is \(25/36\); their posterior means are exactly the rule values in (2). The posterior square-loss identity thus proves a Bayes lower bound \(25/36\) separately at both assignments, hence under every design. Therefore every base singleton risk equals \(5/6\).

For an admissible \(n=2B\), the budget inequalities yield a physical assignment in each selector cell. For the lower bound, draw one atom from (3) and replicate it in all blocks. The normalized target is the base target and the full observation is repeated base data, so the conditional Bayes risk remains \(25/36\) under every physical assignment. For the upper bound, put half the mass on one assignment per selector cell, apply (2) blockwise, and average. Jensen's inequality gives
\[
 \left\{B^{-1}\sum_j(\widehat t_j-t_j)\right\}^2
 \le B^{-1}\sum_j(\widehat t_j-t_j)^2,
\]
and each expected term is at most \(25/36\). Thus every replicated singleton risk is \(5/6\). Taking the stated maxima proves the result. \(\square\)